\section{Frontier Models}
\sectionkicker{FRONTIER}
\begin{wideflow}
{\Large\bfseries 提供条件が具体化した最前線――GPT-6とHaiku 5.5\par}
\smallskip
{\color{SurveyMuted}GPT-6の提供階層とHaiku 5.5の料金区分が具体化した。開始遅延の定義と二段階料金の読み方が実務の分かれ目になる。\par}
\end{wideflow}

\subsection*{有料も無料も、新しい顔ぶれ――GPT-6とHaiku 5.5の提供条件を読み解く}

OpenAIは10月7日、GPT-6と新しい対話インターフェースIntelligent UIの提供を開始した。対象はPlus、Pro、Business、Enterpriseの有料プランで、利用場所はChatタブに限られる。翌10月8日にはFreeとGoプランにも展開された。ここでモデル割り当ての関係が重要になる。有料プランにはSol、FreeとGoにはLunaが割り当てられ、同じGPT-6の看板のもとで提供されるモデルが利用者層によって分かれている。いずれも発行元の提供ページに明記された範囲の事実であり、料金体系そのものはこのページには掲載されていない。\cite{gpt6ui}

この tier 分けは単なる価格差ではない。Chatタブ限定という提供面の指定と、Sol／Lunaというモデル名の対応が同時に示されている点が技術的に意味を持つ。利用者が触れるのは抽象的な「GPT-6」ではなく、プランごとに異なる推論モデルだからだ。提供条件の読み取りでは、モデル名・提供日・利用場所の三点を切り離さずに押さえる必要がある。これらはいずれも提供ページの記述をそのまま受け取ったもので、第三者による検証を経た数値ではない。\cite{gpt6ui}

\begin{table}[htbp]
\centering\footnotesize
\begin{tabularx}{0.98\linewidth}{@{}lX@{}}
\toprule
項目 & 内容 \\
\midrule
GPT-6＋Intelligent UI & 提供開始：10月7日（有料4プラン、Chatタブ）／10月8日（Free・Go）；モデル割当：有料Sol、無料Luna；遅延の定義：回答開始時間（インターリーブ思考・回答方式による）；料金：本ページに記載なし \\
Claude Haiku 5.5 & 提供開始：10月7日；料金：10万token境界の二段階料金、cache-readは\$0.01（10万以下）／\$0.05（10万超）、Sonnet 5.5のcache-read半額（\$0.20→\$0.10）。入力・出力本体の単価や課金単位の別は取得本文に示されていない（この数値だけでの利用費試算は不可、正式料金表の確認が必要）；ベンチマーク：発行元計測表（GDPval 1620、OSWorld2.1 72.4\%、HLE 45.9／57.4、TB4 39.2\%。各評価の対象・条件・尺度は取得本文に示されていない）；system-card全文は未取得 \\
\multicolumn{2}{@{}X@{}}{Ultrafast tier（関連）：gpt-6.1-sol向けResponses APIのservice\_tier、詳細はエージェント基盤の記事で扱う。} \\
\bottomrule
\end{tabularx}
\caption{比較表：フロンティア2件の提供条件と遅延定義}
\end{table}
\noindent \cite{gpt6ui}\cite{haiku55}

GPT-6の発表で最も誤読しやすいのが「44\%高速化」という数字だ。提供ページの記述を正確に読むと、これはGPT-6 Instantと5.6 Instantをweb検索質問のサブセットで比べたとき、回答の開始が44\%早まったという発行元測定である。仕組みは思考と回答を交互に進めるインターリーブ方式によるもので、最終回答が出そろうまでの壁時計時間（wall-time）が44\%縮んだという主張ではない。開始遅延と全体遅延は別の量であり、提供ページは前者だけを測っている。このtime-to-start限定は本記事全体で維持する。\cite{gpt6ui}

この区別は実務に直結する。体感のキレを左右する初動の速さは改善していても、長い回答全体の所要時間が同じ比率で縮む保証はどこにも書かれていない。壁時計時間の短縮を示す数値は提供ページに存在しない。system-cardの数値も列挙されておらず、料金も同ページにない。開始遅延の改善として評価し、全体遅延の主張として扱わないことが正確な読み方だ。\cite{gpt6ui}

\subsection*{Haiku 5.5の二段階料金――10万tokenの境界線}

Anthropicは10月7日、Claude Haiku 5.5を公開した。料金設計の要点は10万tokenを境とする二段階制で、境界以下と以上で単価が分かれる。取得資料からはcache-readの料金区分を確認できる。10万以下が\$0.01、10万超が\$0.05とされる。同時にSonnet 5.5のcache-read料金が\$0.20から\$0.10へ半額化された。ここで分けておきたいのは、入力・出力・cache-readという料金区分の違いだ。取得本文が数値を示すのはcache-readの二区分であり、入力・出力本体の単価や、\$0.01・\$0.05の課金単位の別までは示されていない。課金単位が未確認であるため、この数値だけで実際の利用費を試算することはできない。料金見積もりには正式な料金表の確認が必要となる。自タスクのtoken量が境界のどちら側に寄るかの確認は、正式料金表と併せて行いたい。\cite{haiku55}

Sonnetのcache-read半額化は、Haiku単体の話題を超えて上位モデルの運用費に波及する。キャッシュ読み取りを多用するリトリーバル拡張生成や反復推論では、入力tokenの実効単価が下がる方向に働く。ただし、どのワークロードでどれだけ効くかを示す実測は公開ページになく、各自の利用形態での検証が残る。価格表の読み取りとしては確定情報、効果の見積もりとしては未検証という二層で受け止めるのが正確だ。\cite{haiku55}

性能面では、GDPval 1620、OSWorld2.1で72.4\%、HLEで45.9および57.4、TB4で39.2\%というベンチマーク表が示されている。これらはすべて発行元による測定であり、独立した第三者の再現を経ていない。さらに、取得本文が示すのは評価名と数値のみで、GDPvalが何を、OSWorldが何を、HLEとTB4が何を測るのか、実験条件や尺度は示されていない。表の数値をモデルの優劣の確定と読むのは時期尚早で、あくまで発行元が自社条件で測った申告値として扱うべきだ。評価対象の説明と条件の開示はsystem-card側の取得待ちであり、現時点では不明とする。\cite{haiku55}

Haikuの測定条件の詳細はsystem-card側にしかない可能性があるが、その全文評価やプラットフォーム別の識別子までは取得していない。ベンチ表の行間にある条件の深掘りはsystem-card取得後の課題として残る。現時点では測定条件の留保付きで表を引用するのが限度だ。\cite{haiku55}

両者の確認状況は対照的だ。GPT-6側は提供階層・提供日・モデル割当・開始遅延の定義が確定し、料金とsystem-card数値は未確認である。Haiku側は二段階料金とcache-read区分・ベンチ表の値が確定し、課金単位の別とsystem-card全文・測定条件が未確認である。提供条件の層ではGPT-6が厚く、価格表の層ではHaikuが厚いという違いを保ったまま読み分けることが、誤読を防ぐ近道だ。\cite{gpt6ui}\cite{haiku55}

今週を象徴するのは、2社がそろって利用条件の具体化を進めてきたことだ。GPT-6は提供階層と開始遅延の定義で、Haiku 5.5は二段階料金と申告ベンチ表で、それぞれ開発者の意思決定に直結する材料を置いた。どちらも能力の飛躍をうたう発表ではなく、利用条件の確定に重心がある。ただし価格の確定度は非対称で、Haikuのcache-read区分は数値付きで確定する一方、GPT-6の料金は未確認のまま残る。地味に見えるが、実務への効き目はこうした確定情報のほうが大きい。\cite{gpt6ui}\cite{haiku55}

残る未確認事項を整理する。GPT-6側は料金とsystem-card数値、Haiku側はsystem-card全文とプラットフォーム別識別子が未取得である。ベンチマーク類はすべて独立再現待ちで、本号はいずれも発行元測定として明示した。評価対象の説明がないベンチ表は、追試の出発点としても条件不足であり、system-cardの開示が次の条件になる。\cite{gpt6ui}\cite{haiku55}

日本の読者にとっての要点は三つだ。第一に、利用プランとモデルの対応を確認してから試すこと。第二に、44\%の意味を取り違えて調達判断をしないこと。第三に、自タスクのtoken量がHaikuの10万token境界のどちら側に該当するかを確認し、実際の費用は正式な料金表で課金単位と各料金区分を確認してから試算すること。提供条件の確認は今週確定した材料で進められる一方、費用の金額計算には正式な料金表の確認が要る。料金区分の確認と費用の金額計算は別の作業として扱いたい。\cite{gpt6ui}\cite{haiku55}

\begin{technicalnote}[Note]
本記事の数値はすべて発行元の申告値であり、独立した再現測定ではない。料金表は公開ページの記載を引用し、入力・出力本体の単価や課金単位の別は取得本文に示されていない。
\end{technicalnote}
